%% Generated by tools/gen_error_codes.py from docs/errors/ of the compiler;
%% edit the compiler's pages or tools/error_themes.txt, then rerun it.

\clearpage
\section{Lexical structure}
\label{sec:codes:lexical}

\begin{longtable}{@{}l p{0.78\linewidth} r@{}}
\toprule Code & Message & Page\\ \midrule \endhead
\hyperref[sec:codes:e1001]{E1001} & unterminated string literal \errhole{} & \pageref{sec:codes:e1001}\\
\hyperref[sec:codes:e1002]{E1002} & unterminated comment block \errhole{} & \pageref{sec:codes:e1002}\\
\hyperref[sec:codes:e1004]{E1004} & line break encountered inside a single-line string literal & \pageref{sec:codes:e1004}\\
\hyperref[sec:codes:e1005]{E1005} & escaping a line break is prohibited; use a multiline string instead """ & \pageref{sec:codes:e1005}\\
\hyperref[sec:codes:e1006]{E1006} & escaping a line break is prohibited; use a multiline comment instead /**/ & \pageref{sec:codes:e1006}\\
\hyperref[sec:codes:e1007]{E1007} & escape character found at end of string literal & \pageref{sec:codes:e1007}\\
\hyperref[sec:codes:e1008]{E1008} & escape character found at end of comment block & \pageref{sec:codes:e1008}\\
\hyperref[sec:codes:e1009]{E1009} & invalid numeric value formatted in hexadecimal & \pageref{sec:codes:e1009}\\
\hyperref[sec:codes:e1010]{E1010} & invalid numeric value formatted in decimal & \pageref{sec:codes:e1010}\\
\hyperref[sec:codes:e1011]{E1011} & invalid numeric value formatted in binary & \pageref{sec:codes:e1011}\\
\hyperref[sec:codes:e1012]{E1012} & invalid numeric value formatted in octal & \pageref{sec:codes:e1012}\\
\hyperref[sec:codes:e1013]{E1013} & expected sign information in hexadecimal floating-point literal & \pageref{sec:codes:e1013}\\
\hyperref[sec:codes:e1014]{E1014} & expected sign information in scientific floating-point literal & \pageref{sec:codes:e1014}\\
\hyperref[sec:codes:e1015]{E1015} & \errhole{} is not a valid identifier; use backquotes for special identifiers(e.g. `\errhole{}`) & \pageref{sec:codes:e1015}\\
\hyperref[sec:codes:e1016]{E1016} & binary number is empty & \pageref{sec:codes:e1016}\\
\hyperref[sec:codes:e1017]{E1017} & octal number is empty & \pageref{sec:codes:e1017}\\
\hyperref[sec:codes:e1018]{E1018} & hexadecimal number is empty & \pageref{sec:codes:e1018}\\
\hyperref[sec:codes:e1019]{E1019} & the exponent part of the floating-point number is empty & \pageref{sec:codes:e1019}\\
\hyperref[sec:codes:e1020]{E1020} & the fractional part of the floating-point number is empty & \pageref{sec:codes:e1020}\\
\hyperref[sec:codes:e1021]{E1021} & missing the exponent part in a floating-point number formatted in hexadecimal & \pageref{sec:codes:e1021}\\
\hyperref[sec:codes:e1022]{E1022} & unterminated version block: missing \#end & \pageref{sec:codes:e1022}\\
\hyperref[sec:codes:e1023]{E1023} & \#\errhole{} without matching \#if & \pageref{sec:codes:e1023}\\
\hyperref[sec:codes:e1024]{E1024} & missing version condition: expected identifier or backquoted identifier & \pageref{sec:codes:e1024}\\
\hyperref[sec:codes:e1025]{E1025} & undefined version directive; expecting if, else or end & \pageref{sec:codes:e1025}\\
\bottomrule
\end{longtable}

\errorcode{E1001}{unterminated string literal \errhole{}}
\label{sec:codes:e1001}

Raised during \textbf{lexing}, from \texttt{LexingErrorMessage::\allowbreak{}UNTERMINATED\_\allowbreak{}STRING} (\textit{src/\allowbreak{}ymirc/\allowbreak{}lexing/\allowbreak{}errors.yr}).

\subsubsection*{What it means}

The lexer reached the end of the file while still inside a string literal: the opening quote it reports has no matching closing quote. The quote it was looking for is shown in the message, so a \tokennolst{"} and a \tokennolst{"""} are told apart.

\subsubsection*{Example}

\textit{test\_\allowbreak{}resources/\allowbreak{}lexing/\allowbreak{}test12.yr}:

%% check: skip
\begin{lstlisting}[style=coloredverbatimError]
let x = """abc
\end{lstlisting}

\begin{lstlisting}[style=bashVerb, breaklines=true, escapechar=~]
Error[E1001] : unterminated string literal """
 --> main.yr:(1,9)
 1  ~\lstbox{\textSFxi{}}~ let x = """abc
    ~\lstbox{\textSFv{}}~         ^^^
    ~\lstbox{\textSFii{}}~~\lstbox{\textSFx{}}~~\lstbox{\textSFx{}}~~\lstbox{\textSFx{}}~~\lstbox{\textSFx{}}~~\lstbox{\textSFx{}}~~\lstbox{\textSFx{}}~~\lstbox{\textSFx{}}~
\end{lstlisting}

\subsubsection*{How to fix it}

Close the literal with the same delimiter that opened it. A literal that is meant to span several lines needs the \tokennolst{"""} form, since a plain \tokennolst{"} literal cannot cross a line break (see \hyperref[sec:codes:e1004]{E1004}).

\errorcode{E1002}{unterminated comment block \errhole{}}
\label{sec:codes:e1002}

Raised during \textbf{lexing}, from \texttt{LexingErrorMessage::\allowbreak{}UNTERMINATED\_\allowbreak{}COMMENT} (\textit{src/\allowbreak{}ymirc/\allowbreak{}lexing/\allowbreak{}errors.yr}).

\subsubsection*{What it means}

The lexer reached the end of the file while still inside a comment block. The opening marker it reports was never matched by its closing one.

\subsubsection*{Example}

\textit{test\_\allowbreak{}resources/\allowbreak{}lexing/\allowbreak{}test13.yr}:

%% check: skip
\begin{lstlisting}[style=coloredverbatimError]
fn main () {
    /* unterminated
\end{lstlisting}

\begin{lstlisting}[style=bashVerb, breaklines=true, escapechar=~]
Error[E1002] : unterminated comment block */
 --> main.yr:(2,5)
 2  ~\lstbox{\textSFxi{}}~     /* unterminated
    ~\lstbox{\textSFv{}}~     ^^
    ~\lstbox{\textSFii{}}~~\lstbox{\textSFx{}}~~\lstbox{\textSFx{}}~~\lstbox{\textSFx{}}~~\lstbox{\textSFx{}}~~\lstbox{\textSFx{}}~~\lstbox{\textSFx{}}~~\lstbox{\textSFx{}}~
\end{lstlisting}

\subsubsection*{How to fix it}

Close the comment. \tokennolst{/*} is closed by \tokennolst{*/} and \tokennolst{/**} by \tokennolst{*/} as well; a comment opened inside another one has to be closed too, since block comments nest.

\errorcode{E1004}{line break encountered inside a single-line string literal}
\label{sec:codes:e1004}

Raised during \textbf{lexing}, from \texttt{LexingErrorMessage::\allowbreak{}RETURN\_\allowbreak{}IN\_\allowbreak{}SINGLE\_\allowbreak{}LINE\_\allowbreak{}STRING} (\textit{src/\allowbreak{}ymirc/\allowbreak{}lexing/\allowbreak{}errors.yr}).

\subsubsection*{What it means}

A \tokennolst{"} string literal contains a raw line break. A single-line literal ends at the end of its line, so the closing quote can never be reached.

\subsubsection*{Example}

\textit{test\_\allowbreak{}resources/\allowbreak{}lexing/\allowbreak{}test14.yr}:

%% check: skip
\begin{lstlisting}[style=coloredverbatimError]
let x = "abc
let y = 1;
\end{lstlisting}

\begin{lstlisting}[style=bashVerb, breaklines=true, escapechar=~]
Error[E1004] : line break encountered inside a single-line string literal
 --> main.yr:(1,9)
 1  ~\lstbox{\textSFxi{}}~ let x = "abc

    ~\lstbox{\textSFv{}}~         ^   ^
    ~\lstbox{\textSFii{}}~~\lstbox{\textSFx{}}~~\lstbox{\textSFx{}}~~\lstbox{\textSFx{}}~~\lstbox{\textSFx{}}~~\lstbox{\textSFx{}}~~\lstbox{\textSFx{}}~~\lstbox{\textSFx{}}~
\end{lstlisting}

\subsubsection*{How to fix it}

Either close the literal before the end of the line and concatenate the pieces, or use the multiline form \tokennolst{""" ... """}, which accepts line breaks. \tokennolst{\textbackslash{}n} inside a single-line literal also works when a line break is all that is wanted.

\errorcode{E1005}{escaping a line break is prohibited; use a multiline string instead """}
\label{sec:codes:e1005}

Raised during \textbf{lexing}, from \texttt{LexingErrorMessage::\allowbreak{}ESCAPE\_\allowbreak{}RETURN\_\allowbreak{}PROHIBITED\_\allowbreak{}STRING} (\textit{src/\allowbreak{}ymirc/\allowbreak{}lexing/\allowbreak{}errors.yr}).

\subsubsection*{What it means}

A backslash was used to continue a \tokennolst{"} string literal onto the next line. Ymir does not accept an escaped line break inside a string: a literal either fits on one line or uses the multiline form.

\subsubsection*{Example}

\textit{test\_\allowbreak{}resources/\allowbreak{}lexing/\allowbreak{}test15.yr}:

%% check: skip
\begin{lstlisting}[style=coloredverbatimError]
let x = "abc\
def";
\end{lstlisting}

\begin{lstlisting}[style=bashVerb, breaklines=true, escapechar=~]
Error[E1005] : escaping a line break is prohibited; use a multiline string instead """
 --> main.yr:(1,9)
 1  ~\lstbox{\textSFxi{}}~ let x = "abc\
    ~\lstbox{\textSFv{}}~         ^   ^
    ~\lstbox{\textSFii{}}~~\lstbox{\textSFx{}}~~\lstbox{\textSFx{}}~~\lstbox{\textSFx{}}~~\lstbox{\textSFx{}}~~\lstbox{\textSFx{}}~~\lstbox{\textSFx{}}~~\lstbox{\textSFx{}}~
\end{lstlisting}

\subsubsection*{How to fix it}

Write the literal as \tokennolst{""" ... """}, which spans lines without any escaping, or put \tokennolst{\textbackslash{}n} in a single-line literal where the break is wanted.

\errorcode{E1006}{escaping a line break is prohibited; use a multiline comment instead /**/}
\label{sec:codes:e1006}

Raised during \textbf{lexing}, from \texttt{LexingErrorMessage::\allowbreak{}ESCAPE\_\allowbreak{}RETURN\_\allowbreak{}PROHIBITED\_\allowbreak{}COMMENT} (\textit{src/\allowbreak{}ymirc/\allowbreak{}lexing/\allowbreak{}errors.yr}).

\subsubsection*{What it means}

A backslash was used to continue a single-line comment onto the next line. Escaping a line break has no meaning in a comment.

\subsubsection*{Example}

\textit{test\_\allowbreak{}resources/\allowbreak{}lexing/\allowbreak{}test16.yr}:

%% check: skip
\begin{lstlisting}[style=coloredverbatimError]
// abc\
def
\end{lstlisting}

\begin{lstlisting}[style=bashVerb, breaklines=true, escapechar=~]
Error[E1006] : escaping a line break is prohibited; use a multiline comment instead /**/
 --> main.yr:(1,1)
 1  ~\lstbox{\textSFxi{}}~ // abc\
    ~\lstbox{\textSFv{}}~ ^^    ^
    ~\lstbox{\textSFii{}}~~\lstbox{\textSFx{}}~~\lstbox{\textSFx{}}~~\lstbox{\textSFx{}}~~\lstbox{\textSFx{}}~~\lstbox{\textSFx{}}~~\lstbox{\textSFx{}}~~\lstbox{\textSFx{}}~
\end{lstlisting}

\subsubsection*{How to fix it}

Use a block comment (\tokennolst{/* ... */}) for a comment that spans several lines, or repeat \tokennolst{//} on each line.

\errorcode{E1007}{escape character found at end of string literal}
\label{sec:codes:e1007}

Raised during \textbf{lexing}, from \texttt{LexingErrorMessage::\allowbreak{}ESCAPE\_\allowbreak{}AT\_\allowbreak{}END\_\allowbreak{}OF\_\allowbreak{}STRING} (\textit{src/\allowbreak{}ymirc/\allowbreak{}lexing/\allowbreak{}errors.yr}).

\subsubsection*{What it means}

The string literal ends on a backslash: the escape sequence it starts has nothing to escape, because the closing quote is the next character.

\subsubsection*{Example}

\textit{test\_\allowbreak{}resources/\allowbreak{}lexing/\allowbreak{}test17.yr}:

%% check: skip
\begin{lstlisting}[style=coloredverbatimError]
let x = "abc\
\end{lstlisting}

\begin{lstlisting}[style=bashVerb, breaklines=true, escapechar=~]
Error[E1007] : escape character found at end of string literal
 --> main.yr:(1,9)
 1  ~\lstbox{\textSFxi{}}~ let x = "abc\
    ~\lstbox{\textSFv{}}~         ^   ^
    ~\lstbox{\textSFii{}}~~\lstbox{\textSFx{}}~~\lstbox{\textSFx{}}~~\lstbox{\textSFx{}}~~\lstbox{\textSFx{}}~~\lstbox{\textSFx{}}~~\lstbox{\textSFx{}}~~\lstbox{\textSFx{}}~
\end{lstlisting}

\subsubsection*{How to fix it}

Write \tokennolst{\textbackslash{}\textbackslash{}} if a literal backslash is wanted, or complete the escape sequence with the character it applies to.

\errorcode{E1008}{escape character found at end of comment block}
\label{sec:codes:e1008}

Raised during \textbf{lexing}, from \texttt{LexingErrorMessage::\allowbreak{}ESCAPE\_\allowbreak{}AT\_\allowbreak{}END\_\allowbreak{}OF\_\allowbreak{}COMMENT} (\textit{src/\allowbreak{}ymirc/\allowbreak{}lexing/\allowbreak{}errors.yr}).

\subsubsection*{What it means}

The comment block ends on a backslash, so the escape sequence it starts has nothing to escape.

\subsubsection*{Example}

\textit{test\_\allowbreak{}resources/\allowbreak{}lexing/\allowbreak{}test18.yr}:

%% check: skip
\begin{lstlisting}[style=coloredverbatimError]
/* abc\
\end{lstlisting}

\begin{lstlisting}[style=bashVerb, breaklines=true, escapechar=~]
Error[E1008] : escape character found at end of comment block
 --> main.yr:(1,1)
 1  ~\lstbox{\textSFxi{}}~ /* abc\
    ~\lstbox{\textSFv{}}~ ^^    ^
    ~\lstbox{\textSFii{}}~~\lstbox{\textSFx{}}~~\lstbox{\textSFx{}}~~\lstbox{\textSFx{}}~~\lstbox{\textSFx{}}~~\lstbox{\textSFx{}}~~\lstbox{\textSFx{}}~~\lstbox{\textSFx{}}~
\end{lstlisting}

\subsubsection*{How to fix it}

Remove the trailing backslash, or double it if a literal backslash is wanted.

\errorcode{E1009}{invalid numeric value formatted in hexadecimal}
\label{sec:codes:e1009}

Raised during \textbf{lexing}, from \texttt{LexingErrorMessage::\allowbreak{}INVALID\_\allowbreak{}HEX\_\allowbreak{}NUMBER} (\textit{src/\allowbreak{}ymirc/\allowbreak{}lexing/\allowbreak{}errors.yr}).

\subsubsection*{What it means}

A literal introduced by \tokennolst{0x} contains a character that is not a hexadecimal digit. Hexadecimal digits are \tokennolst{0} to \tokennolst{9}, \tokennolst{a} to \tokennolst{f} and \tokennolst{A} to \tokennolst{F}, plus \tokennolst{\_} as a separator.

\subsubsection*{Example}

\textit{test\_\allowbreak{}resources/\allowbreak{}lexing/\allowbreak{}test23.yr}:

%% check: skip
\begin{lstlisting}[style=coloredverbatimError]
let x = 0xZ1;
\end{lstlisting}

\begin{lstlisting}[style=bashVerb, breaklines=true, escapechar=~]
Error[E1009] : invalid numeric value formatted in hexadecimal
 --> main.yr:(1,9)
 1  ~\lstbox{\textSFxi{}}~ let x = 0xZ1;
    ~\lstbox{\textSFv{}}~         ^^^^
    ~\lstbox{\textSFxi{}}~ Note :
    ~\lstbox{\textSFxi{}}~  --> main.yr:(1,11)
    ~\lstbox{\textSFxi{}}~  1  ~\lstbox{\textSFxi{}}~ let x = 0xZ1;
    ~\lstbox{\textSFxi{}}~     ~\lstbox{\textSFv{}}~           ^
    ~\lstbox{\textSFxi{}}~     ~\lstbox{\textSFii{}}~~\lstbox{\textSFx{}}~~\lstbox{\textSFx{}}~~\lstbox{\textSFx{}}~~\lstbox{\textSFx{}}~~\lstbox{\textSFx{}}~~\lstbox{\textSFx{}}~~\lstbox{\textSFx{}}~
    ~\lstbox{\textSFii{}}~~\lstbox{\textSFx{}}~~\lstbox{\textSFx{}}~~\lstbox{\textSFx{}}~~\lstbox{\textSFx{}}~~\lstbox{\textSFx{}}~~\lstbox{\textSFvii{}}~~\lstbox{\textSFx{}}~
\end{lstlisting}

\subsubsection*{How to fix it}

Remove the offending character. If it was meant to be a type suffix, check its spelling: a suffix the lexer does not know is read as part of the number.

\errorcode{E1010}{invalid numeric value formatted in decimal}
\label{sec:codes:e1010}

Raised during \textbf{lexing}, from \texttt{LexingErrorMessage::\allowbreak{}INVALID\_\allowbreak{}DEC\_\allowbreak{}NUMBER} (\textit{src/\allowbreak{}ymirc/\allowbreak{}lexing/\allowbreak{}errors.yr}).

\subsubsection*{What it means}

A decimal literal contains a character that is neither a digit nor a valid suffix.

\subsubsection*{Example}

\textit{test\_\allowbreak{}resources/\allowbreak{}lexing/\allowbreak{}test25.yr}:

%% check: skip
\begin{lstlisting}[style=coloredverbatimError]
let x = 12a4;
\end{lstlisting}

\begin{lstlisting}[style=bashVerb, breaklines=true, escapechar=~]
Error[E1010] : invalid numeric value formatted in decimal
 --> main.yr:(1,9)
 1  ~\lstbox{\textSFxi{}}~ let x = 12a4;
    ~\lstbox{\textSFv{}}~         ^^^^
    ~\lstbox{\textSFxi{}}~ Note :
    ~\lstbox{\textSFxi{}}~  --> main.yr:(1,11)
    ~\lstbox{\textSFxi{}}~  1  ~\lstbox{\textSFxi{}}~ let x = 12a4;
    ~\lstbox{\textSFxi{}}~     ~\lstbox{\textSFv{}}~           ^
    ~\lstbox{\textSFxi{}}~     ~\lstbox{\textSFii{}}~~\lstbox{\textSFx{}}~~\lstbox{\textSFx{}}~~\lstbox{\textSFx{}}~~\lstbox{\textSFx{}}~~\lstbox{\textSFx{}}~~\lstbox{\textSFx{}}~~\lstbox{\textSFx{}}~
    ~\lstbox{\textSFii{}}~~\lstbox{\textSFx{}}~~\lstbox{\textSFx{}}~~\lstbox{\textSFx{}}~~\lstbox{\textSFx{}}~~\lstbox{\textSFx{}}~~\lstbox{\textSFvii{}}~~\lstbox{\textSFx{}}~
\end{lstlisting}

\subsubsection*{How to fix it}

Remove the offending character, or separate the number from what follows it. \tokennolst{\_} is accepted inside a number as a digit separator (\tokennolst{1\_000\_000}).

\errorcode{E1011}{invalid numeric value formatted in binary}
\label{sec:codes:e1011}

Raised during \textbf{lexing}, from \texttt{LexingErrorMessage::\allowbreak{}INVALID\_\allowbreak{}BIN\_\allowbreak{}NUMBER} (\textit{src/\allowbreak{}ymirc/\allowbreak{}lexing/\allowbreak{}errors.yr}).

\subsubsection*{What it means}

A literal introduced by \tokennolst{0b} contains a character other than \tokennolst{0} or \tokennolst{1}.

\subsubsection*{Example}

\textit{test\_\allowbreak{}resources/\allowbreak{}lexing/\allowbreak{}test19.yr}:

%% check: skip
\begin{lstlisting}[style=coloredverbatimError]
let x = 0b1021;
\end{lstlisting}

\begin{lstlisting}[style=bashVerb, breaklines=true, escapechar=~]
Error[E1011] : invalid numeric value formatted in binary
 --> main.yr:(1,9)
 1  ~\lstbox{\textSFxi{}}~ let x = 0b1021;
    ~\lstbox{\textSFv{}}~         ^^^^^^
    ~\lstbox{\textSFxi{}}~ Note :
    ~\lstbox{\textSFxi{}}~  --> main.yr:(1,13)
    ~\lstbox{\textSFxi{}}~  1  ~\lstbox{\textSFxi{}}~ let x = 0b1021;
    ~\lstbox{\textSFxi{}}~     ~\lstbox{\textSFv{}}~             ^
    ~\lstbox{\textSFxi{}}~     ~\lstbox{\textSFii{}}~~\lstbox{\textSFx{}}~~\lstbox{\textSFx{}}~~\lstbox{\textSFx{}}~~\lstbox{\textSFx{}}~~\lstbox{\textSFx{}}~~\lstbox{\textSFx{}}~~\lstbox{\textSFx{}}~
    ~\lstbox{\textSFii{}}~~\lstbox{\textSFx{}}~~\lstbox{\textSFx{}}~~\lstbox{\textSFx{}}~~\lstbox{\textSFx{}}~~\lstbox{\textSFx{}}~~\lstbox{\textSFvii{}}~~\lstbox{\textSFx{}}~
\end{lstlisting}

\subsubsection*{How to fix it}

Use only \tokennolst{0} and \tokennolst{1} in a binary literal, with \tokennolst{\_} as an optional separator. Write the value in another base if binary is not what was wanted.

\errorcode{E1012}{invalid numeric value formatted in octal}
\label{sec:codes:e1012}

Raised during \textbf{lexing}, from \texttt{LexingErrorMessage::\allowbreak{}INVALID\_\allowbreak{}OCT\_\allowbreak{}NUMBER} (\textit{src/\allowbreak{}ymirc/\allowbreak{}lexing/\allowbreak{}errors.yr}).

\subsubsection*{What it means}

A literal introduced by \tokennolst{0o} contains a digit outside \tokennolst{0} to \tokennolst{7}.

\subsubsection*{Example}

\textit{test\_\allowbreak{}resources/\allowbreak{}lexing/\allowbreak{}test21.yr}:

%% check: skip
\begin{lstlisting}[style=coloredverbatimError]
let x = 0o1892;
\end{lstlisting}

\begin{lstlisting}[style=bashVerb, breaklines=true, escapechar=~]
Error[E1012] : invalid numeric value formatted in octal
 --> main.yr:(1,9)
 1  ~\lstbox{\textSFxi{}}~ let x = 0o1892;
    ~\lstbox{\textSFv{}}~         ^^^^^^
    ~\lstbox{\textSFxi{}}~ Note :
    ~\lstbox{\textSFxi{}}~  --> main.yr:(1,12)
    ~\lstbox{\textSFxi{}}~  1  ~\lstbox{\textSFxi{}}~ let x = 0o1892;
    ~\lstbox{\textSFxi{}}~     ~\lstbox{\textSFv{}}~            ^
    ~\lstbox{\textSFxi{}}~     ~\lstbox{\textSFii{}}~~\lstbox{\textSFx{}}~~\lstbox{\textSFx{}}~~\lstbox{\textSFx{}}~~\lstbox{\textSFx{}}~~\lstbox{\textSFx{}}~~\lstbox{\textSFx{}}~~\lstbox{\textSFx{}}~
    ~\lstbox{\textSFii{}}~~\lstbox{\textSFx{}}~~\lstbox{\textSFx{}}~~\lstbox{\textSFx{}}~~\lstbox{\textSFx{}}~~\lstbox{\textSFx{}}~~\lstbox{\textSFvii{}}~~\lstbox{\textSFx{}}~
\end{lstlisting}

\subsubsection*{How to fix it}

Use only octal digits. A leading zero does not make a number octal in Ymir -- the \tokennolst{0o} prefix does -- so \tokennolst{09} is a decimal literal, not an octal one.

\errorcode{E1013}{expected sign information in hexadecimal floating-point literal}
\label{sec:codes:e1013}

Raised during \textbf{lexing}, from \texttt{LexingErrorMessage::\allowbreak{}MISSING\_\allowbreak{}HEX\_\allowbreak{}FLOAT\_\allowbreak{}SIGN} (\textit{src/\allowbreak{}ymirc/\allowbreak{}lexing/\allowbreak{}errors.yr}).

\subsubsection*{What it means}

A hexadecimal floating-point literal has a \tokennolst{p} exponent marker that is not followed by a sign. The exponent of a hexadecimal float is written with an explicit \tokennolst{+} or \tokennolst{-}.

\subsubsection*{Example}

\textit{test\_\allowbreak{}resources/\allowbreak{}lexing/\allowbreak{}test28.yr}:

%% check: skip
\begin{lstlisting}[style=coloredverbatimError]
let x = 0x1.8p;
\end{lstlisting}

\begin{lstlisting}[style=bashVerb, breaklines=true, escapechar=~]
Error[E1013] : expected sign information in hexadecimal floating-point literal
 --> main.yr:(1,9)
 1  ~\lstbox{\textSFxi{}}~ let x = 0x1.8p;
    ~\lstbox{\textSFv{}}~         ^^^
    ~\lstbox{\textSFii{}}~~\lstbox{\textSFx{}}~~\lstbox{\textSFx{}}~~\lstbox{\textSFx{}}~~\lstbox{\textSFx{}}~~\lstbox{\textSFx{}}~~\lstbox{\textSFx{}}~~\lstbox{\textSFx{}}~
\end{lstlisting}

\subsubsection*{How to fix it}

Write the sign: \tokennolst{0x1.8p+3} rather than \tokennolst{0x1.8p3}.

\errorcode{E1014}{expected sign information in scientific floating-point literal}
\label{sec:codes:e1014}

Raised during \textbf{lexing}, from \texttt{LexingErrorMessage::\allowbreak{}MISSING\_\allowbreak{}SCI\_\allowbreak{}FLOAT\_\allowbreak{}SIGN} (\textit{src/\allowbreak{}ymirc/\allowbreak{}lexing/\allowbreak{}errors.yr}).

\subsubsection*{What it means}

A floating-point literal in scientific notation has an \tokennolst{e} exponent marker that is not followed by a sign.

\subsubsection*{Example}

\textit{test\_\allowbreak{}resources/\allowbreak{}lexing/\allowbreak{}test27.yr}:

%% check: skip
\begin{lstlisting}[style=coloredverbatimError]
let x = 1.0e;
\end{lstlisting}

\begin{lstlisting}[style=bashVerb, breaklines=true, escapechar=~]
Error[E1014] : expected sign information in scientific floating-point literal
 --> main.yr:(1,9)
 1  ~\lstbox{\textSFxi{}}~ let x = 1.0e;
    ~\lstbox{\textSFv{}}~         ^
    ~\lstbox{\textSFii{}}~~\lstbox{\textSFx{}}~~\lstbox{\textSFx{}}~~\lstbox{\textSFx{}}~~\lstbox{\textSFx{}}~~\lstbox{\textSFx{}}~~\lstbox{\textSFx{}}~~\lstbox{\textSFx{}}~
\end{lstlisting}

\subsubsection*{How to fix it}

Write the sign: \tokennolst{1.5e+10} rather than \tokennolst{1.5e10}.

\errorcode{E1015}{\errhole{} is not a valid identifier; use backquotes for special identifiers(e.g. `\errhole{}`)}
\label{sec:codes:e1015}

Raised during \textbf{lexing}, from \texttt{LexingErrorMessage::\allowbreak{}INVALID\_\allowbreak{}IDENTIFIER} (\textit{src/\allowbreak{}ymirc/\allowbreak{}lexing/\allowbreak{}errors.yr}).

\subsubsection*{What it means}

The name reported is not a valid identifier: it either starts with a character identifiers cannot start with, or it is a reserved keyword. Ymir can still name such a symbol, but the name has to be quoted.

\subsubsection*{Example}

\textit{test\_\allowbreak{}resources/\allowbreak{}lexing/\allowbreak{}test30.yr}:

%% check: skip
\begin{lstlisting}[style=coloredverbatimError]
let __ = 1;
\end{lstlisting}

\begin{lstlisting}[style=bashVerb, breaklines=true, escapechar=~]
Error[E1015] : __ is not a valid identifier; use backquotes for special identifiers(e.g. `__`)
 --> main.yr:(1,5)
 1  ~\lstbox{\textSFxi{}}~ let __ = 1;
    ~\lstbox{\textSFv{}}~     ^^
    ~\lstbox{\textSFii{}}~~\lstbox{\textSFx{}}~~\lstbox{\textSFx{}}~~\lstbox{\textSFx{}}~~\lstbox{\textSFx{}}~~\lstbox{\textSFx{}}~~\lstbox{\textSFx{}}~~\lstbox{\textSFx{}}~
\end{lstlisting}

\subsubsection*{How to fix it}

Rename the symbol, or write it between backquotes -- \tokennolst{`if`} -- which makes any text usable as an identifier. This is how a symbol coming from another language whose name collides with an Ymir keyword is referred to.

\errorcode{E1016}{binary number is empty}
\label{sec:codes:e1016}

Raised during \textbf{lexing}, from \texttt{LexingErrorMessage::\allowbreak{}EMPTY\_\allowbreak{}BIN\_\allowbreak{}NUMBER} (\textit{src/\allowbreak{}ymirc/\allowbreak{}lexing/\allowbreak{}errors.yr}).

\subsubsection*{What it means}

A \tokennolst{0b} prefix is not followed by any binary digit, so the literal has no value.

\subsubsection*{Example}

\textit{test\_\allowbreak{}resources/\allowbreak{}lexing/\allowbreak{}test20.yr}:

%% check: skip
\begin{lstlisting}[style=coloredverbatimError]
let x = 0b_;
\end{lstlisting}

\begin{lstlisting}[style=bashVerb, breaklines=true, escapechar=~]
Error[E1016] : binary number is empty
 --> main.yr:(1,9)
 1  ~\lstbox{\textSFxi{}}~ let x = 0b_;
    ~\lstbox{\textSFv{}}~         ^^^
    ~\lstbox{\textSFii{}}~~\lstbox{\textSFx{}}~~\lstbox{\textSFx{}}~~\lstbox{\textSFx{}}~~\lstbox{\textSFx{}}~~\lstbox{\textSFx{}}~~\lstbox{\textSFx{}}~~\lstbox{\textSFx{}}~
\end{lstlisting}

\subsubsection*{How to fix it}

Write the digits after the prefix, or remove the prefix if a plain \tokennolst{0} was meant.

\errorcode{E1017}{octal number is empty}
\label{sec:codes:e1017}

Raised during \textbf{lexing}, from \texttt{LexingErrorMessage::\allowbreak{}EMPTY\_\allowbreak{}OCT\_\allowbreak{}NUMBER} (\textit{src/\allowbreak{}ymirc/\allowbreak{}lexing/\allowbreak{}errors.yr}).

\subsubsection*{What it means}

A \tokennolst{0o} prefix is not followed by any octal digit, so the literal has no value.

\subsubsection*{Example}

\textit{test\_\allowbreak{}resources/\allowbreak{}lexing/\allowbreak{}test22.yr}:

%% check: skip
\begin{lstlisting}[style=coloredverbatimError]
let x = 0o_;
\end{lstlisting}

\begin{lstlisting}[style=bashVerb, breaklines=true, escapechar=~]
Error[E1017] : octal number is empty
 --> main.yr:(1,9)
 1  ~\lstbox{\textSFxi{}}~ let x = 0o_;
    ~\lstbox{\textSFv{}}~         ^^^
    ~\lstbox{\textSFii{}}~~\lstbox{\textSFx{}}~~\lstbox{\textSFx{}}~~\lstbox{\textSFx{}}~~\lstbox{\textSFx{}}~~\lstbox{\textSFx{}}~~\lstbox{\textSFx{}}~~\lstbox{\textSFx{}}~
\end{lstlisting}

\subsubsection*{How to fix it}

Write the digits after the prefix, or remove the prefix if a plain \tokennolst{0} was meant.

\errorcode{E1018}{hexadecimal number is empty}
\label{sec:codes:e1018}

Raised during \textbf{lexing}, from \texttt{LexingErrorMessage::\allowbreak{}EMPTY\_\allowbreak{}HEX\_\allowbreak{}NUMBER} (\textit{src/\allowbreak{}ymirc/\allowbreak{}lexing/\allowbreak{}errors.yr}).

\subsubsection*{What it means}

A \tokennolst{0x} prefix is not followed by any hexadecimal digit, so the literal has no value.

\subsubsection*{Example}

\textit{test\_\allowbreak{}resources/\allowbreak{}lexing/\allowbreak{}test24.yr}:

%% check: skip
\begin{lstlisting}[style=coloredverbatimError]
let x = 0x_;
\end{lstlisting}

\begin{lstlisting}[style=bashVerb, breaklines=true, escapechar=~]
Error[E1018] : hexadecimal number is empty
 --> main.yr:(1,9)
 1  ~\lstbox{\textSFxi{}}~ let x = 0x_;
    ~\lstbox{\textSFv{}}~         ^^^
    ~\lstbox{\textSFii{}}~~\lstbox{\textSFx{}}~~\lstbox{\textSFx{}}~~\lstbox{\textSFx{}}~~\lstbox{\textSFx{}}~~\lstbox{\textSFx{}}~~\lstbox{\textSFx{}}~~\lstbox{\textSFx{}}~
\end{lstlisting}

\subsubsection*{How to fix it}

Write the digits after the prefix, or remove the prefix if a plain \tokennolst{0} was meant.

\errorcode{E1019}{the exponent part of the floating-point number is empty}
\label{sec:codes:e1019}

Raised during \textbf{lexing}, from \texttt{LexingErrorMessage::\allowbreak{}EMPTY\_\allowbreak{}EXP\_\allowbreak{}FLOAT\_\allowbreak{}NUMBER} (\textit{src/\allowbreak{}ymirc/\allowbreak{}lexing/\allowbreak{}errors.yr}).

\subsubsection*{What it means}

The exponent marker of a floating-point literal is followed by a sign but by no digit, so the exponent has no value.

\subsubsection*{Example}

\textit{test\_\allowbreak{}resources/\allowbreak{}lexing/\allowbreak{}test5.yr}:

%% check: skip
\begin{lstlisting}[style=coloredverbatimError]
1.e+
\end{lstlisting}

\begin{lstlisting}[style=bashVerb, breaklines=true, escapechar=~]
Error[E1019] : the exponent part of the floating-point number is empty
 --> main.yr:(1,1)
 1  ~\lstbox{\textSFxi{}}~ 1.e+
    ~\lstbox{\textSFv{}}~ ^
    ~\lstbox{\textSFii{}}~~\lstbox{\textSFx{}}~~\lstbox{\textSFx{}}~~\lstbox{\textSFx{}}~~\lstbox{\textSFx{}}~~\lstbox{\textSFx{}}~~\lstbox{\textSFx{}}~~\lstbox{\textSFx{}}~
\end{lstlisting}

\subsubsection*{How to fix it}

Write the exponent digits, as in \tokennolst{1.0e+10}, or drop the exponent part entirely.

\errorcode{E1020}{the fractional part of the floating-point number is empty}
\label{sec:codes:e1020}

Raised during \textbf{lexing}, from \texttt{LexingErrorMessage::\allowbreak{}EMPTY\_\allowbreak{}DEC\_\allowbreak{}FLOAT\_\allowbreak{}NUMBER} (\textit{src/\allowbreak{}ymirc/\allowbreak{}lexing/\allowbreak{}errors.yr}).

\subsubsection*{What it means}

A floating-point literal has a decimal point with no digit after it.

\subsubsection*{Example}

\textit{test\_\allowbreak{}resources/\allowbreak{}lexing/\allowbreak{}test26.yr}:

%% check: skip
\begin{lstlisting}[style=coloredverbatimError]
let x = 1._;
\end{lstlisting}

\begin{lstlisting}[style=bashVerb, breaklines=true, escapechar=~]
Error[E1020] : the fractional part of the floating-point number is empty
 --> main.yr:(1,9)
 1  ~\lstbox{\textSFxi{}}~ let x = 1._;
    ~\lstbox{\textSFv{}}~         ^
    ~\lstbox{\textSFxi{}}~ Note :
    ~\lstbox{\textSFxi{}}~  --> main.yr:(1,11)
    ~\lstbox{\textSFxi{}}~  1  ~\lstbox{\textSFxi{}}~ let x = 1._;
    ~\lstbox{\textSFxi{}}~     ~\lstbox{\textSFv{}}~           ^
    ~\lstbox{\textSFxi{}}~     ~\lstbox{\textSFii{}}~~\lstbox{\textSFx{}}~~\lstbox{\textSFx{}}~~\lstbox{\textSFx{}}~~\lstbox{\textSFx{}}~~\lstbox{\textSFx{}}~~\lstbox{\textSFx{}}~~\lstbox{\textSFx{}}~
    ~\lstbox{\textSFii{}}~~\lstbox{\textSFx{}}~~\lstbox{\textSFx{}}~~\lstbox{\textSFx{}}~~\lstbox{\textSFx{}}~~\lstbox{\textSFx{}}~~\lstbox{\textSFvii{}}~~\lstbox{\textSFx{}}~
\end{lstlisting}

\subsubsection*{How to fix it}

Write the fractional digits (\tokennolst{1.0}), or drop the point if an integer was meant. A trailing point is not a valid way to write a float in Ymir.

\errorcode{E1021}{missing the exponent part in a floating-point number formatted in hexadecimal}
\label{sec:codes:e1021}

Raised during \textbf{lexing}, from \texttt{LexingErrorMessage::\allowbreak{}MISSING\_\allowbreak{}P\_\allowbreak{}IN\_\allowbreak{}HEX\_\allowbreak{}FLOAT} (\textit{src/\allowbreak{}ymirc/\allowbreak{}lexing/\allowbreak{}errors.yr}).

\subsubsection*{What it means}

A hexadecimal literal contains a decimal point, which makes it a floating-point literal, but it has no \tokennolst{p} exponent part. A hexadecimal float always carries its binary exponent explicitly.

\subsubsection*{Example}

\textit{test\_\allowbreak{}resources/\allowbreak{}lexing/\allowbreak{}test29.yr}:

%% check: skip
\begin{lstlisting}[style=coloredverbatimError]
let x = 0x1.8;
\end{lstlisting}

\begin{lstlisting}[style=bashVerb, breaklines=true, escapechar=~]
Error[E1021] : missing the exponent part in a floating-point number formatted in hexadecimal
 --> main.yr:(1,9)
 1  ~\lstbox{\textSFxi{}}~ let x = 0x1.8;
    ~\lstbox{\textSFv{}}~         ^^^
    ~\lstbox{\textSFxi{}}~ Note :
    ~\lstbox{\textSFxi{}}~  --> main.yr:(1,12)
    ~\lstbox{\textSFxi{}}~  1  ~\lstbox{\textSFxi{}}~ let x = 0x1.8;
    ~\lstbox{\textSFxi{}}~     ~\lstbox{\textSFv{}}~            ^
    ~\lstbox{\textSFxi{}}~     ~\lstbox{\textSFii{}}~~\lstbox{\textSFx{}}~~\lstbox{\textSFx{}}~~\lstbox{\textSFx{}}~~\lstbox{\textSFx{}}~~\lstbox{\textSFx{}}~~\lstbox{\textSFx{}}~~\lstbox{\textSFx{}}~
    ~\lstbox{\textSFii{}}~~\lstbox{\textSFx{}}~~\lstbox{\textSFx{}}~~\lstbox{\textSFx{}}~~\lstbox{\textSFx{}}~~\lstbox{\textSFx{}}~~\lstbox{\textSFvii{}}~~\lstbox{\textSFx{}}~
\end{lstlisting}

\subsubsection*{How to fix it}

Add the exponent: \tokennolst{0x1.8p+0}. Use a decimal literal instead if the exponent notation is not what was wanted.

\errorcode{E1022}{unterminated version block: missing \#end}
\label{sec:codes:e1022}

Raised during \textbf{lexing}, from \texttt{LexingErrorMessage::\allowbreak{}UNTERMINATED\_\allowbreak{}VERSION\_\allowbreak{}BLOCK} (\textit{src/\allowbreak{}ymirc/\allowbreak{}lexing/\allowbreak{}errors.yr}).

\subsubsection*{What it means}

A \tokennolst{\#if} conditional compilation block was opened and the file ended before the matching \tokennolst{\#end} was found.

\subsubsection*{Example}

\textit{test\_\allowbreak{}resources/\allowbreak{}lexing/\allowbreak{}test9.yr}:

%% check: skip
\begin{lstlisting}[style=coloredverbatimError]
#if FOO
\end{lstlisting}

\begin{lstlisting}[style=bashVerb, breaklines=true, escapechar=~]
Error[E1022] : unterminated version block: missing #end
 --> main.yr:(1,2)
 1  ~\lstbox{\textSFxi{}}~ #if FOO
    ~\lstbox{\textSFv{}}~  ^^
    ~\lstbox{\textSFii{}}~~\lstbox{\textSFx{}}~~\lstbox{\textSFx{}}~~\lstbox{\textSFx{}}~~\lstbox{\textSFx{}}~~\lstbox{\textSFx{}}~~\lstbox{\textSFx{}}~~\lstbox{\textSFx{}}~
\end{lstlisting}

\subsubsection*{How to fix it}

Close the block with \tokennolst{\#end}. Every \tokennolst{\#if} needs one, including the ones whose \tokennolst{\#else} branch is empty.

\errorcode{E1023}{\#\errhole{} without matching \#if}
\label{sec:codes:e1023}

Raised during \textbf{lexing}, from \texttt{LexingErrorMessage::\allowbreak{}UNMATCHED\_\allowbreak{}VERSION\_\allowbreak{}END} (\textit{src/\allowbreak{}ymirc/\allowbreak{}lexing/\allowbreak{}errors.yr}).

\subsubsection*{What it means}

A \tokennolst{\#else} or \tokennolst{\#end} directive was found with no \tokennolst{\#if} open at that point -- either the \tokennolst{\#if} is missing, or an earlier \tokennolst{\#end} already closed it.

\subsubsection*{Example}

\textit{test\_\allowbreak{}resources/\allowbreak{}lexing/\allowbreak{}test8.yr}:

%% check: skip
\begin{lstlisting}[style=coloredverbatimError]
#else
\end{lstlisting}

\begin{lstlisting}[style=bashVerb, breaklines=true, escapechar=~]
Error[E1023] : #else without matching #if
 --> main.yr:(1,2)
 1  ~\lstbox{\textSFxi{}}~ #else
    ~\lstbox{\textSFv{}}~  ^^^^
    ~\lstbox{\textSFii{}}~~\lstbox{\textSFx{}}~~\lstbox{\textSFx{}}~~\lstbox{\textSFx{}}~~\lstbox{\textSFx{}}~~\lstbox{\textSFx{}}~~\lstbox{\textSFx{}}~~\lstbox{\textSFx{}}~
\end{lstlisting}

\subsubsection*{How to fix it}

Add the missing \tokennolst{\#if}, or remove the stray directive. Check the nesting of the enclosing blocks: an \tokennolst{\#end} written one level too early closes the outer block and leaves the next directive unmatched.

\errorcode{E1024}{missing version condition: expected identifier or backquoted identifier}
\label{sec:codes:e1024}

Raised during \textbf{lexing}, from \texttt{LexingErrorMessage::\allowbreak{}MISSING\_\allowbreak{}VERSION\_\allowbreak{}CONDITION} (\textit{src/\allowbreak{}ymirc/\allowbreak{}lexing/\allowbreak{}errors.yr}).

\subsubsection*{What it means}

A \tokennolst{\#if} directive is not followed by the version identifier it should test.

\subsubsection*{Example}

\textit{test\_\allowbreak{}resources/\allowbreak{}lexing/\allowbreak{}test7.yr}:

%% check: skip
\begin{lstlisting}[style=coloredverbatimError]
#if
foo
#end
\end{lstlisting}

\begin{lstlisting}[style=bashVerb, breaklines=true, escapechar=~]
Error[E1024] : missing version condition: expected identifier or backquoted identifier
 --> main.yr:(1,2)
 1  ~\lstbox{\textSFxi{}}~ #if
    ~\lstbox{\textSFv{}}~  ^^
    ~\lstbox{\textSFii{}}~~\lstbox{\textSFx{}}~~\lstbox{\textSFx{}}~~\lstbox{\textSFx{}}~~\lstbox{\textSFx{}}~~\lstbox{\textSFx{}}~~\lstbox{\textSFx{}}~~\lstbox{\textSFx{}}~
\end{lstlisting}

\subsubsection*{How to fix it}

Name the version to test after the directive, as in \tokennolst{\#if linux}. An identifier that is not a plain word has to be backquoted.

\errorcode{E1025}{undefined version directive; expecting if, else or end}
\label{sec:codes:e1025}

Raised during \textbf{lexing}, from \texttt{LexingErrorMessage::\allowbreak{}UNDEFINED\_\allowbreak{}VERSION\_\allowbreak{}DIRECTIVE} (\textit{src/\allowbreak{}ymirc/\allowbreak{}lexing/\allowbreak{}errors.yr}).

\subsubsection*{What it means}

The word following \tokennolst{\#} is not one of the directives the lexer knows. The only conditional compilation directives are \tokennolst{\#if}, \tokennolst{\#else} and \tokennolst{\#end}.

\subsubsection*{Example}

\textit{test\_\allowbreak{}resources/\allowbreak{}lexing/\allowbreak{}test10.yr}:

%% check: skip
\begin{lstlisting}[style=coloredverbatimError]
#foo
\end{lstlisting}

\begin{lstlisting}[style=bashVerb, breaklines=true, escapechar=~]
Error[E1025] : undefined version directive; expecting if, else or end
 --> main.yr:(1,1)
 1  ~\lstbox{\textSFxi{}}~ #foo
    ~\lstbox{\textSFv{}}~ ^
    ~\lstbox{\textSFxi{}}~ Note :
    ~\lstbox{\textSFxi{}}~  --> main.yr:(1,2)
    ~\lstbox{\textSFxi{}}~  1  ~\lstbox{\textSFxi{}}~ #foo
    ~\lstbox{\textSFxi{}}~     ~\lstbox{\textSFv{}}~  ^^^
    ~\lstbox{\textSFxi{}}~     ~\lstbox{\textSFii{}}~~\lstbox{\textSFx{}}~~\lstbox{\textSFx{}}~~\lstbox{\textSFx{}}~~\lstbox{\textSFx{}}~~\lstbox{\textSFx{}}~~\lstbox{\textSFx{}}~~\lstbox{\textSFx{}}~
    ~\lstbox{\textSFii{}}~~\lstbox{\textSFx{}}~~\lstbox{\textSFx{}}~~\lstbox{\textSFx{}}~~\lstbox{\textSFx{}}~~\lstbox{\textSFx{}}~~\lstbox{\textSFvii{}}~~\lstbox{\textSFx{}}~
\end{lstlisting}

\subsubsection*{How to fix it}

Correct the spelling of the directive, or remove the \tokennolst{\#} if the line was not meant to be a directive at all.
